\documentclass[11pt,a4paper]{article}

\usepackage[a4paper,margin=2.5cm]{geometry}
\usepackage[utf8]{inputenc}
\usepackage[T1]{fontenc}
\usepackage{lmodern}
\usepackage{titlesec}
\usepackage{fancyhdr}
\usepackage{enumitem}
\usepackage{array}
\usepackage{longtable}

\titleformat{\section}{\normalfont\Large\bfseries}{\thesection}{0.6em}{}
\titlespacing*{\section}{0pt}{1.3em}{0.5em}
\titleformat{\subsection}{\normalfont\large\bfseries}{\thesubsection}{0.6em}{}
\titlespacing*{\subsection}{0pt}{1em}{0.4em}

\pagestyle{fancy}
\fancyhf{}
\renewcommand{\headrulewidth}{0pt}
\fancyfoot[C]{\small YANFLÈ --- Cahier des charges --- \thepage}

\setlength{\parindent}{0pt}
\setlength{\parskip}{0.45em}

\newcommand{\regle}[1]{\vspace{0.3em}\noindent\textit{#1}\vspace{0.3em}}

\begin{document}

\begin{center}
  {\Huge\bfseries YANFLÈ}\\[6pt]
  {\large Cahier des charges}\\[4pt]
  {\itshape Tech Impact --- Orange Summer Challenge 2026}
\end{center}

\vspace{1em}

Ce document décrit le projet YANFLÈ tel qu'il est envisagé dans son ensemble : un système de vidéosurveillance capable de comprendre les situations qu'il filme et de prévenir un danger avant qu'il ne se produise, et non un simple détecteur de mouvement. Il ne décrit pas ce qui a déjà été codé (voir pour cela le document d'architecture technique du prototype), mais ce que le projet doit accomplir dans sa version complète.

\section{Présentation du contexte}

\subsection{Le problème}
Les caméras de surveillance actuelles sont passives : elles enregistrent, mais n'analysent rien de ce qu'elles filment. Le détecteur de mouvement classique ne distingue pas un événement anodin d'une situation dangereuse. La caméra ne sert donc qu'à constater les faits après coup, une fois le vol, le cambriolage ou l'agression déjà commis.

\subsection{Objectif du projet}
Faire passer la vidéosurveillance d'un rôle de témoin (qui ne sert qu'à prouver qu'un acte a eu lieu) à un rôle de vigie, capable d'aider à l'empêcher, en s'appuyant sur les caméras déjà installées.

\subsection{Les acteurs}
\begin{itemize}[leftmargin=1.4em, itemsep=2pt]
  \item \textbf{Maîtrise d'ouvrage (MOA)} : les futurs sites clients (responsables de magasins, d'entreprises, de bâtiments) qui expriment le besoin de sécurité, dans le cadre de l'Orange Summer Challenge 2026.
  \item \textbf{Maîtrise d'œuvre (MOE)} : l'équipe Tech Impact, qui conçoit et réalise le système.
  \item \textbf{Utilisateurs finaux} : les responsables de site et les opérateurs de sécurité qui reçoivent et traitent les alertes.
\end{itemize}

\section{Périmètre et besoins fonctionnels}

\subsection{Périmètre inclus}
\begin{itemize}[leftmargin=1.4em, itemsep=2pt]
  \item Analyse en temps réel des flux vidéo de plusieurs caméras, réparties sur plusieurs bâtiments.
  \item Reconnaissance de comportements à risque (et pas seulement de mouvement).
  \item Reconnaissance faciale utilisée uniquement pour gérer les autorisations de présence par lieu (pas d'identification à d'autres fins).
  \item Alerte graduée : simple notification pour un signal faible, appel vocal automatisé pour une situation critique.
  \item Une interface d'administration multi-bâtiment et multi-caméra pour les responsables de site.
  \item Un assistant conversationnel permettant d'interroger l'état d'un bâtiment (présences, caméras, alertes en cours).
\end{itemize}

\subsection{Périmètre exclu}
\begin{itemize}[leftmargin=1.4em, itemsep=2pt]
  \item Toute décision ou action automatique sans validation humaine : le système alerte, il ne décide jamais seul.
  \item L'identification de personnes à des fins judiciaires ou l'utilisation des données biométriques en dehors de la gestion d'autorisation interne.
  \item Le stockage des flux vidéo sur un serveur cloud tiers : le traitement reste local, sur un serveur du client.
  \item La génération de descriptions ou d'alertes non vérifiables : le système ne rapporte que des faits observés, jamais une supposition présentée comme un fait.
\end{itemize}

\subsection{Besoins fonctionnels principaux (User Stories)}
\begin{itemize}[leftmargin=1.4em, itemsep=4pt]
  \item En tant qu'opérateur de sécurité, je veux être alerté automatiquement d'un comportement suspect, afin de pouvoir intervenir avant qu'un incident ne se produise.
  \item En tant qu'opérateur de sécurité, je veux recevoir un appel vocal en cas de situation critique, afin d'être certain que l'alerte a bien été prise en compte même si je ne suis pas devant un écran.
  \item En tant que responsable de site, je veux définir les lieux où chaque personne reconnue est autorisée à se trouver, afin d'être informé d'une présence anormale.
  \item En tant que responsable de site, je veux gérer plusieurs bâtiments et caméras depuis une seule interface, afin de superviser plusieurs sites avec une seule équipe.
  \item En tant que responsable de site, je veux interroger un assistant sur l'état de mon bâtiment, afin d'obtenir une information rapide sans consulter chaque caméra une par une.
\end{itemize}

\section{Contraintes techniques et opérationnelles}
\begin{itemize}[leftmargin=1.4em, itemsep=2pt]
  \item Le système s'appuie sur des caméras IP standard, sans matériel propriétaire imposé.
  \item Le traitement (analyse vidéo, reconnaissance, alertes) s'exécute sur un serveur local équipé GPU, sans dépendance à un service cloud externe pour fonctionner.
  \item Toute communication impliquant un flux vidéo (caméra vers serveur, serveur vers interface) est chiffrée (HTTPS/TLS).
  \item Le traitement d'images de personnes est soumis au RGPD : durée de conservation limitée, droit à l'image, finalité strictement limitée à la sécurité du site.
\end{itemize}

\section{Besoins non fonctionnels}
\begin{itemize}[leftmargin=1.4em, itemsep=2pt]
  \item \textbf{Temps réel} : une situation à risque doit être signalée en quelques secondes, pas après coup.
  \item \textbf{Disponibilité} : fonctionnement continu, avec reprise automatique après une coupure (réseau, redémarrage serveur).
  \item \textbf{Fiabilité} : le taux de fausses alertes doit rester assez bas pour ne pas décrédibiliser le système auprès des opérateurs.
  \item \textbf{Sécurité d'accès} : chaque caméra et chaque bâtiment est protégé par un code d'accès qui lui est propre.
  \item \textbf{Scalabilité} : il doit être possible d'ajouter un bâtiment ou une caméra sans réinstallation complète du système.
\end{itemize}

\section{Priorisation (MoSCoW)}

\begin{longtable}{@{}p{0.14\textwidth}p{0.82\textwidth}@{}}
\textbf{Must have} & Analyse comportementale temps réel multi-caméra ; alerte graduée jusqu'à l'appel vocal ; reconnaissance faciale limitée à la gestion d'autorisation par lieu ; traitement local (pas de cloud) ; aucune alerte basée sur une information inventée. \\[6pt]
\textbf{Should have} & Assistant conversationnel pour interroger l'état d'un bâtiment ; historique consultable des alertes passées. \\[6pt]
\textbf{Could have} & Application mobile dédiée aux opérateurs ; tableaux de bord statistiques sur les comportements détectés. \\[6pt]
\textbf{Won't have} & Identification de personnes à des fins judiciaires ; action automatique sans validation humaine ; hébergement des flux vidéo sur un cloud public. \\
\end{longtable}

\section{Livrables, calendrier et suivi}

\subsection{Livrables}
Code source du système, documentation technique, manuel d'utilisation pour les responsables de site, modèles entraînés pour la reconnaissance de comportements à risque.

\subsection{Calendrier}
Le projet avance par étapes successives, chacune livrant une version utilisable :
\begin{enumerate}[leftmargin=1.4em, itemsep=2pt]
  \item Preuve de concept mono-caméra (détection et alertes de base).
  \item Système multi-caméra et multi-bâtiment, avec gestion des autorisations.
  \item Entraînement d'un modèle dédié à la reconnaissance des signaux précurseurs (avant l'acte, pas seulement après).
  \item Escalade automatisée par appel vocal selon la gravité.
  \item Déploiement pilote sur un site réel.
\end{enumerate}

\subsection{Suivi}
Développement en cycles courts, avec démonstration régulière de l'avancement à l'équipe Tech Impact et ajustement du périmètre si nécessaire.

\vspace{1em}
\noindent\textit{Ce cahier des charges décrit l'ambition complète du projet. La version actuellement développée en est une première étape (voir le document d'architecture technique du prototype pour le détail de ce qui est réellement fait à ce jour).}

\end{document}
